\section{Solutions for Chapter~\ref{ch:defects}}

\paragraph{Exercise~\ref{ex:defect-projector}.}
\phantomsection
\label{sol:ex:defect-projector}
On the orthogonal summand $M_i\otimes\overline{M_i}$, the
projector candidate has eigenvalue
\[
 p_i=\frac1{\mathcal D^2}\sum_a
       \frac{S_{a0}}{S_{00}}\frac{S_{ai}}{S_{0i}}
 =\frac{\delta_{0i}}{\mathcal D^2S_{00}S_{0i}}
 =\delta_{0i}.
\]
Here the middle equality is orthogonality of the real $S$
columns, and the last uses $S_{00}=1/\mathcal D$ when $i=0$.
Therefore $\Pi_{\rm inv}$ squares
to itself and is Hermitian on the orthogonal sector
decomposition. Its vacuum eigenvalue before normalization
is $\sum_a d_a^2=\mathcal D^2$, which explains the
denominator. Dividing by $\mathcal D$ would leave the
vacuum eigenvalue $\mathcal D$ and would fail idempotence.



\paragraph{Exercise~\ref{ex:defect-twined}.}
\phantomsection
\label{sol:ex:defect-twined}
On the scalar blocks, $\alpha=1$ and the two possible
$\rho$ eigenvalues solve $x^2=1+3x$.
On the two-dimensional block,
$\alpha^3=1$, $\rho^2=1$, and conjugation by the swap
matrix sends $\alpha$ to $\alpha^{-1}$. Moreover
$1+\alpha+\alpha^2=0$, so
$\rho+\alpha\rho+\alpha^2\rho=0$ there. These checks
verify the generating relations and hence all products.
To make the noncommutative products explicit, put
$r_i=\alpha^i\rho$ with indices modulo three. Then
\[
 \alpha^j r_i=r_{i+j},\qquad
 r_i\alpha^j=r_{i-j},\qquad
 r_i r_j=\alpha^{i-j}\rho^2
        =\alpha^{i-j}+r_0+r_1+r_2.
\]
The middle equality in the last formula follows by moving
$\alpha^j$ past $\rho$ with $\rho\alpha^j=\alpha^{-j}\rho$.
The traces of $\alpha$ on the three multiplicity blocks
are $1,1,-1$; those of $\rho$ are $d,-d^{-1},0$.
Multiplying by the respective characters and the unchanged
antiholomorphic factor gives the two twined traces.
For example, on $\C^2\otimes a$ the first-factor operator is
$\widehat\alpha|_{\C^2}\otimes I_a$, so its graded trace
is $(\omega+\omega^2)\operatorname{Tr}_a q^{L_0-c/24}
=-\chi_a$. The matrix acts on the two-copy multiplicity
space, not on a two-dimensional approximation to the infinite
module $a$.



\paragraph{Exercise~\ref{ex:defect-alpha}.}
\phantomsection
\label{sol:ex:defect-alpha}
The forgetful table gives
$R\alpha=c_0\oplus c_1\oplus c_2$.
The three summands have lowest weights $1$, $1/3$, and $2/3$,
respectively. Thus the ground comes only from $c_1$.
Its restriction multiplicity is $n_{\alpha,c_1}=1$, whereas
the ordinary dimension of its chiral top is three.
In the first-factor twisted Hilbert space the antiholomorphic
$E_8$ vacuum has multiplicity one, so the lowest full level
has
\[
 (h,\bar h)=\left(\frac13,0\right),\qquad
 \Delta=s_{\rm spin}=\frac13,\qquad
 d_{\rm states}=1\cdot3\cdot1=3.
\]
There is no second contribution with the chiralities exchanged
for this specified first-factor twist.
A twisted state sits at a defect endpoint or in a
defect boundary condition; the ordinary untwisted locality
requirement of integer spin does not apply to it unchanged.



\paragraph{Exercise~\ref{ex:defect-between}.}
\phantomsection
\label{sol:ex:defect-between}
Put $x=\alpha^i\rho$ and $y=\alpha^j\rho$.
Duality reverses the order, and the fusion relation then gives
$x^*=\rho\alpha^{-i}=\alpha^i\rho=x$.
An oriented boundary-changing field from $x$ to $y$ is counted
by $x^*y$, not by deleting the dual in a general formula.
For $i\ne j$,
$(\alpha^i\rho)^*(\alpha^j\rho)
=\alpha^{i-j}+\rho+\alpha\rho+\alpha^2\rho$.
Indeed, $x^*y=\alpha^{i-j}\rho^2$, and multiplying the
three noninvertible summands of $\rho^2$ by $\alpha^{i-j}$
only permutes them. The nontrivial invertible term restricts to
$c_0+c_1+c_2$, contributing three states of weight $1/3$.
Each noninvertible contributes another three. The total is
$3+3\cdot3=12$, compared with nine in the self-boundary
case, where the invertible term is the vacuum boundary.
Equivalently, the character $\chi_{c_1}$ occurs four times
between distinct noninvertible boundaries and three times
in a self-annulus; each copy has three top states.
The distinct-boundary annulus has no identity state:
the coefficient of the tensor unit in $x^*y$ is zero when
$x\ne y$. These are boundary weights $h=1/3$ in a single
chiral open-channel trace, rather than bulk scalar dimensions
$2h=2/3$.



\paragraph{Exercise~\ref{ex:defect-categories}.}
\phantomsection
\label{sol:ex:defect-categories}
Six counts the simples of the Haagerup boundary fusion
category; twelve counts those of its braided Drinfeld center;
thirty-six counts the simples of the relevant bimodule defect
category equivalent, up to the stated opposite convention,
to a product of the boundary category with its opposite.
The center is braided; the boundary and bimodule categories
are used as fusion categories. A Lagrangian algebra has
only the trivial local module category, although its all-module
category has six simples. Replacing all modules by local
ones would discard precisely the nonlocal boundary sectors.

